\documentclass{article}
\usepackage{iftex}
\ifXeTeX
  \usepackage{fontspec}
  \usepackage{xeCJK}
  \usepackage[haranoaji]{zxjafont}
  \XeTeXlinebreaklocale ""
\fi
\ifLuaTeX
  \usepackage[haranoaji]{luatexja-preset}
\fi
\usepackage[a4paper, margin=2cm]{geometry}
\usepackage{graphicx}
\newdimen\charwd
{\tt\global\setbox0=\hbox{x}\global\charwd=\wd0}
\frenchspacing
\usepackage{fancyhdr}
\pagestyle{fancy}
\renewcommand{\headrulewidth}{0pt}
\fancyhf{}
\fancyhead[R]{\rm src2tex-go version 2.23}
% plain TeX compatibility
\makeatletter
\providecommand{\eqalign}[1]{\vcenter{\openup1\jot\m@th
  \ialign{\strut\hfil$\displaystyle{##}$&$\displaystyle{{}##}$\hfil\crcr#1\crcr}}}
\makeatother
\fancyfoot[R]{\rm\hfill samples{\tt /}sqrt{\tt\_}mat.red\qquad page \thepage}
\begin{document}

\ifx\sevenrm\undefined
  \font\sevenrm=cmr7 scaled \magstep0
\fi

\newread\MyStyle
\openin\MyStyle=src2latex.s2t
\ifeof\MyStyle
  \closein\MyStyle
\else
  \input src2latex.s2t
  \closein\MyStyle
\fi

\def\mc{\relax}
\def\gt{\relax}
\def\sc{\scshape}

\tt\mc 

\noindent
\rm\mc {\tt\%}\kern1\charwd  \bf sqrt\_mat.red \rm\mc 

\noindent
\mbox{}\tt\mc \hfill

\noindent
\mbox{}\rm\mc {\tt\%}\kern1\charwd  {}For a given $2\times2$ matrix $\,A\ge0\,$, we shall prove that
 $$
 \sqrt{A}={1\over\sqrt{2\sqrt{ac-b^2}+a+c}}
 \pmatrix{
 \sqrt{ac-b^2}+a & b \cr
 b & \sqrt{ac-b^2}+c \cr}
 $$ \rm\mc 

\noindent
\mbox{}{\tt\%}

\noindent
\mbox{}\tt\mc \hfill

\noindent
\mbox{}on{\tt\mc \kern1\charwd}d\kern-.455em d\kern-.045em i\kern-.455em i\kern-.045em v\kern-.455em v\kern-.045em ;

\noindent
\mbox{}\hfill

\noindent
\mbox{}\rm\mc {\tt\%}\kern1\charwd  {}For a given matrix
 $
 A=\pmatrix{ a & b \cr b & c \cr}
 $
 ,\ we shall find a matrix
 $
 X=\pmatrix{ x & y \cr y & z \cr}
 $
 satisfying $A=X^2$. \rm\mc 

\noindent
\mbox{}{\tt\%}

\noindent
\mbox{}\tt\mc \hfill

\noindent
\mbox{}mat{\tt\_\kern.141em}a:=mat((a,b),(b,c));

\noindent
\mbox{}mat{\tt\_\kern.141em}x:=mat((x,y),(y,z));

\noindent
\mbox{}\hfill

\noindent
\mbox{}\rm\mc {\tt\%}\kern1\charwd  {}The equation
 $\,A=X^2\,$ is translated into algebraic equations
 $$
 x^2+y^2=a,\ y(x+z)=b,\ y^2+z^2=c \leqno(1)
 $$ \rm\mc 

\noindent
\mbox{}{\tt\%}

\noindent
\mbox{}\tt\mc \hfill

\noindent
\mbox{}mat{\tt\_\kern.141em}xx:=mat{\tt\_\kern.141em}x{\tt *}mat{\tt\_\kern.141em}x;

\noindent
\mbox{}\hfill

\noindent
\mbox{}\rm\mc {\tt\%}\kern1\charwd  {}(1) is equivalent to
 $$
 x=\sqrt{a-y^2},
	\ y\bigl(\sqrt{a-y^2}+\sqrt{c-y^2}\,\bigr)=b,
	\ z=\sqrt{c-y^2} \leqno(2)
 $$
 We put
 $$
 f(y)=\sqrt{a-y^2},
 \ g(y)=y\bigl(\sqrt{a-y^2}+\sqrt{c-y^2}\,\bigr)-b,
 \ h(y)=\sqrt{c-y^2} \leqno(3)
 $$ \rm\mc 

\noindent
\mbox{}{\tt\%}

\noindent
\mbox{}\tt\mc \hfill

\noindent
\mbox{}func{\tt\_\kern.141em}f:=rhs(first(solve(mat{\tt\_\kern.141em}xx(1,1)=a,x)));

\noindent
\mbox{}func{\tt\_\kern.141em}h:=rhs(first(solve(mat{\tt\_\kern.141em}xx(2,2)=c,z)));

\noindent
\mbox{}func{\tt\_\kern.141em}g:=sub({\tt\char'173}x=func{\tt\_\kern.141em}f,z=func{\tt\_\kern.141em}h{\tt\char'175},mat{\tt\_\kern.141em}xx(1,2)){\tt -}b;

\noindent
\mbox{}\hfill

\noindent
\mbox{}\rm\mc {\tt\%}\kern1\charwd  {}We define a function $\,g_1(w)\,$ by
 $$
 g_1(w)=g(y)\;{\sqrt{a-y^2}-\sqrt{c-y^2}\over y}\;\Bigg|_{y=1/w}
	=-b\sqrt{aw^2-1}+b\sqrt{cw^2-1}+a-c \leqno(4)
 $$ \rm\mc 

\noindent
\mbox{}{\tt\%}

\noindent
\mbox{}\tt\mc \hfill

\noindent
\mbox{}func{\tt\_\kern.141em}g1:=sub(y=1{\tt /}w,func{\tt\_\kern.141em}g{\tt *}(sqrt(a{\tt -}y{\tt\char'136}2){\tt -}sqrt(c{\tt -}y{\tt\char'136}2)){\tt /}y);

\noindent
\mbox{}\hfill

\noindent
\mbox{}\rm\mc {\tt\%}\kern1\charwd  {}We put
 $$
 g_2(v)
	=g_1(w)\Big|_{w=\sqrt{(v^2+1)/a}}
	=b\sqrt{-a+cv^2+c\over a}-bv+a-c \leqno(5)
 $$ \rm\mc 

\noindent
\mbox{}{\tt\%}

\noindent
\mbox{}\tt\mc \hfill

\noindent
\mbox{}func{\tt\_\kern.141em}g2:=sub(w=sqrt((v{\tt\char'136}2+1){\tt /}a),func{\tt\_\kern.141em}g1);

\noindent
\mbox{}\hfill

\noindent
\mbox{}\rm\mc {\tt\%}\kern1\charwd  {}Since
 $$
 \eqalign{
 g_2(v)&=0\cr
 g_2(v)+bv-a+c&=bv-a+c\cr
 (g_2(v)+bv-a+c)^2&=(bv-a+c)^2\cr
 {b^2\over a\ }\;(-a+cv^2+c)&=(bv-a+c)^2\cr
 }
 $$
 we put
 $$
 g_3(v)={b^2\over a\ }\;(-a+cv^2+c)-(bv-a+c)^2 \leqno(6)
 $$
 We can solve the equation $\,g_3(v)=0\,$ and get two solutions
 $$
 v_1={a+\sqrt{ac-b^2}\over b}\ ,\quad v_2={a-\sqrt{ac-b^2}\over b}
 $$
 Here we note that since
 $$
 g_2(v_2)=\sqrt{-2c\sqrt{ac-b^2}+ac-b^2+c^2}+\sqrt{ac-b^2}-c
 $$
 $\,g_2(v_2)\,$ is not always equal to 0. So, we can reject $\,v_2\,$. \rm\mc 

\noindent
\mbox{}{\tt\%}

\noindent
\mbox{}\tt\mc \hfill

\noindent
\mbox{}func{\tt\_\kern.141em}g3:=(func{\tt\_\kern.141em}g2+b{\tt *}v{\tt -}a+c){\tt\char'136}2{\tt -}(b{\tt *}v{\tt -}a+c){\tt\char'136}2;

\noindent
\mbox{}solutions:=solve(func{\tt\_\kern.141em}g3,v);

\noindent
\mbox{}v1:=rhs(first(solutions));

\noindent
\mbox{}v2:=rhs(second(solutions));

\noindent
\mbox{}sub(v=v1,func{\tt\_\kern.141em}g2);

\noindent
\mbox{}sub(v=v2,func{\tt\_\kern.141em}g2);

\noindent
\mbox{}\hfill

\noindent
\mbox{}\rm\mc {\tt\%}\kern1\charwd  {}Since
 $$
 y=\sqrt{a\over v^2+1}
 $$
 we have
 $$
 y_1={b\over \sqrt{2\sqrt{ac-b^2}+a+c}}
 $$ \rm\mc 

\noindent
\mbox{}{\tt\%}

\noindent
\mbox{}\tt\mc \hfill

\noindent
\mbox{}y1:=sqrt(a{\tt /}(v1{\tt\char'136}2+1));

\noindent
\mbox{}\hfill

\noindent
\mbox{}\rm\mc {\tt\%}\kern1\charwd  {}Hence, we obtain
 $$
 x_1=f(y_1)\ ,\quad z_1=h(y_1)
 $$
 this gives
 $$
 X_1
 =\pmatrix{x_1 & y_1 \cr
	y_1 & z_1 \cr}
 ={1\over d_1}\pmatrix{ \sqrt{ac-b^2}+a & b \cr
	b & \sqrt{ac-b^2}+c \cr}
 $$
 where $\,d_1=\sqrt{2\sqrt{ac-b^2}+a+c}\,$.
 Direct computation shows $\,X_1^2=A\,$. \rm\mc 

\noindent
\mbox{}{\tt\%}

\noindent
\mbox{}\tt\mc \hfill

\noindent
\mbox{}x1:=sub(y=y1,func{\tt\_\kern.141em}f);

\noindent
\mbox{}z1:=sub(y=y1,func{\tt\_\kern.141em}h);

\noindent
\mbox{}mat{\tt\_\kern.141em}x1:=mat((x1,y1),(y1,z1));

\noindent
\mbox{}d1:=sqrt(2{\tt *}sqrt(a{\tt *}c{\tt\mc \kern1\charwd}{\tt -}{\tt\mc \kern1\charwd}b{\tt *}{\tt *}2){\tt\mc \kern1\charwd}+{\tt\mc \kern1\charwd}a{\tt\mc \kern1\charwd}+{\tt\mc \kern1\charwd}c);

\noindent
\mbox{}d1{\tt *}mat{\tt\_\kern.141em}x1;

\noindent
\mbox{}mat{\tt\_\kern.141em}x1{\tt *}mat{\tt\_\kern.141em}x1;

\noindent
\mbox{}\hfill

\noindent
\mbox{}e\kern-.445em e\kern-.055em n\kern-.46em n\kern-.04em d\kern-.455em d\kern-.045em ;

\noindent
\mbox{}

\rm\mc

\end{document}
